\section{The Vision of God}

Scripture sets the promise against a wall. Moses on the mountain asks, ``Shew me thy glory,'' and is answered: ``Thou canst not see my face: for man shall not see me, and live.'' He is set in a hole of the rock, covered by the hand of God while the glory passes, and allowed to see only what is behind.\endnote{Ex 33:18, 20--23, Douay--Rheims (Challoner).} The Apostle, praising the King of kings, says that he ``only hath immortality and inhabiteth light inaccessible: whom no man hath seen, nor can see.''\endnote{1 Tim 6:16, Douay--Rheims (Challoner).} And yet the same Lord who covered Moses with his hand says on the mountain of the Beatitudes that the clean of heart shall see God, and the same Apostle writes: ``We see now through a glass in a dark manner: but then face to face. Now I know in part: but then I shall know even as I am known.''\endnote{1 Cor 13:12, Douay--Rheims (Challoner).}

\subsection{The glory of God is a living man}

Among the first Fathers to treat both sides of this at length was Irenaeus of Lyons, writing against those who said that the God the prophets saw could not be the invisible Father. He takes up the hardest text directly. That no man shall see God and live is true, he says, of God's greatness, and he answers it from God's love:

\begin{quote}
But in respect to His greatness, and His wonderful glory, ``no man shall see God and live,'' for the Father is incomprehensible; but in regard to His love, and kindness, and as to His infinite power, even this He grants to those who love Him, that is, to see God, which thing the prophets did also predict. ``For those things that are impossible with men, are possible with God.'' For man does not see God by his own powers; but when He pleases He is seen by men, by whom He wills, and when He wills, and as He wills.\endnote{Irenaeus, \emph{Against Heresies} IV.20.5, trans.\ Alexander Roberts and W.~H. Rambaut, \emph{The Ante-Nicene Fathers}, vol.~1 (Buffalo: Christian Literature Publishing Company, 1887).}
\end{quote}

Then he says what the seeing does to those who see, comparing them to people who stand in light:

\begin{quote}
For as those who see the light are within the light, and partake of its brilliancy; even so, those who see God are in God, and receive of His splendour. But [His] splendour vivifies them; those, therefore, who see God, do receive life.\endnote{Irenaeus, \emph{Against Heresies} IV.20.5, ANF 1.}
\end{quote}

From this he draws a conclusion that answers the words spoken to Moses. Exodus says that no man shall see God and live; Irenaeus writes: ``Men therefore shall see God, that they may live, being made immortal by that sight, and attaining even unto God.''\endnote{Irenaeus, \emph{Against Heresies} IV.20.6, ANF 1.} A few lines later he defines that life by its object:

\begin{quote}
For the glory of God is a living man; and the life of man consists in beholding God.\endnote{Irenaeus, \emph{Against Heresies} IV.20.7, ANF 1.}
\end{quote}

Irenaeus adds that the Word reveals God by degrees, ``preserving at the same time the invisibility of the Father, lest man should at any time become a despiser of God, and that he should always possess something towards which he might advance.''\endnote{Irenaeus, \emph{Against Heresies} IV.20.7, ANF 1.}

\subsection{What the Church confesses}

What Irenaeus taught as a Father the Church confessed as dogma. In 1336 Benedict XII, by the apostolic constitution \emph{Benedictus Deus}, defined that the souls of the saints, once nothing remains in them to be purified, see God even before they take up their bodies again and before the general judgment. The constitution's teaching on the moment after death and on the judgment that follows is set out in \emph{What Death Does Not Finish}. Here only its words about the seeing are needed, taken one phrase at a time:

\begin{quote}
\latin{\ldots\ viderunt et vident divinam essentiam visione intuitiva et etiam faciali, nulla mediante creatura in ratione obiecti visi se habente, sed divina essentia immediate se nude, clare et aperte eis ostendente\ldots}\endnote{Benedict XII, apostolic constitution \emph{Benedictus Deus} (29 January 1336), in H.~Denzinger, \emph{Enchiridion symbolorum}, n.~1000 (Latin).}
\end{quote}

The Catechism renders the heart of it: these souls ``have seen and do see the divine essence with an intuitive vision, and even face to face, without the mediation of any creature.''\endnote{\emph{Catechism of the Catholic Church} 1023, quoting \emph{Benedictus Deus} (DS 1000).}

\emph{\latin{Visione intuitiva}}: by an intuitive vision. The blessed do not reason their way to God from his works, as we do now, climbing from the creature to the Creator; they see him directly.

\emph{\latin{Et etiam faciali}}: and even face to face. The word is Saint Paul's, taken from ``then face to face.'' Augustine had already explained that ``By `the face' of God we are to understand His manifestation, and not a part of the body''.\endnote{Augustine, \emph{The City of God} XXII.29, trans.\ Marcus Dods, \emph{The Works of Aurelius Augustine}, vol.~2 (Edinburgh: T.\,\&\,T.~Clark, 1871).}

\emph{\latin{Nulla mediante creatura in ratione obiecti visi se habente}}: no creature standing between, in the role of the thing seen. What the blessed see is God himself, not an image of him, nor a glory less than God, nor a messenger.

\emph{\latin{Sed divina essentia immediate se nude, clare et aperte eis ostendente}}: but the divine essence showing itself to them immediately, plainly, clearly and openly. In this last clause the subject changes: it is no longer the souls who see but the divine essence, which shows itself to them. The vision is God's gift of himself, and the soul receives it. The Catechism gives the reason: ``Because of his transcendence, God cannot be seen as he is, unless he himself opens up his mystery to man's immediate contemplation and gives him the capacity for it.''\endnote{\emph{Catechism of the Catholic Church} 1028.}

A century later the Council of Florence confessed the same vision and added two words from the epistle of John: the purified see clearly God himself, three and one, as he is: \latin{sicuti est}. In the same sentence it defined that one sees more perfectly than another, \latin{pro meritorum tamen diversitate}, according to the diversity of merits.\endnote{Council of Florence, bull \emph{Laetentur caeli} (6 July 1439), in Denzinger, \emph{Enchiridion symbolorum}, n.~1305 (Latin).}

\subsection{The light of glory}

Thomas Aquinas asks how a creature can see what exceeds every creature, in the question of the \emph{Summa} on how God is known by us. He states the difficulty in an image of his own: God is in himself supremely knowable, ``But what is supremely knowable in itself, may not be knowable to a particular intellect, on account of the excess of the intelligible object above the intellect; as, for example, the sun, which is supremely visible, cannot be seen by the bat by reason of its excess of light.''\endnote{Thomas Aquinas, \emph{Summa theologiae} I, q.~12, a.~1, trans.\ Fathers of the English Dominican Province (New York: Benziger Brothers; Project Gutenberg eBook 17611).} The created mind stands to God as the bat to the sun: by its own nature it cannot look at him. And so Thomas concludes, as Irenaeus had, that the seeing must be given: ``the created intellect cannot see the essence of God, unless God by His grace unites Himself to the created intellect, as an object made intelligible to it.''\endnote{Aquinas, \emph{Summa theologiae} I, q.~12, a.~4, trans.\ English Dominican Province.}

If the seeing must be given, something must change on the creature's side. For this Thomas cites a verse of the psalms, ``In thy light we shall see light,'' and gives an argument:\endnote{Ps 35:10, Douay--Rheims (Challoner); quoted by Aquinas, \emph{Summa theologiae} I, q.~12, a.~5, \emph{sed contra}.}

\begin{quote}
Everything which is raised up to what exceeds its nature, must be prepared by some disposition above its nature; as, for example, if air is to receive the form of fire, it must be prepared by some disposition for such a form. But when any created intellect sees the essence of God, the essence of God itself becomes the intelligible form of the intellect. Hence it is necessary that some supernatural disposition should be added to the intellect in order that it may be raised up to such a great and sublime height. \ldots\ By this light the blessed are made ``deiform''---i.e.\ like to God, according to the saying: ``When He shall appear we shall be like to Him, and [Vulg.: `because'] we shall see Him as He is.''\endnote{Aquinas, \emph{Summa theologiae} I, q.~12, a.~5, trans.\ English Dominican Province.}
\end{quote}

This is the \latin{lumen gloriae}, the light of glory. It strengthens the mind as a habit strengthens a power, and Thomas explains in what sense it can be called a medium:

\begin{quote}
This light is required to see the divine essence, not as a similitude in which God is seen, but as a perfection of the intellect, strengthening it to see God. Therefore it may be said that this light is to be described not as a medium in which God is seen, but as one by which He is seen; and such a medium does not take away the immediate vision of God.\endnote{Aquinas, \emph{Summa theologiae} I, q.~12, a.~5, reply to obj.~2, trans.\ English Dominican Province.}
\end{quote}

Thomas's account agrees with the definition of 1336: no creature stands between the soul and God as the thing seen, because the only created thing involved is the soul itself, made able to see. The Council of Vienne in 1312 condemned the contrary error, that the soul \latin{non indiget lumine gloriae, ipsam elevante ad Deum videndum}, does not need the light of glory raising it to see God.\endnote{Council of Vienne, constitution \emph{Ad nostrum qui} (1312), error 5, in Denzinger, \emph{Enchiridion symbolorum}, n.~895 (Latin).} Such a light, he goes on, ``cannot be natural to a creature unless the creature has a divine nature; which is impossible.''\endnote{Aquinas, \emph{Summa theologiae} I, q.~12, a.~5, reply to obj.~3, trans.\ English Dominican Province.} The vision is therefore the end to which grace brings us, not the growth of a power we were born with.

\subsection{Seeing without exhausting}

Thomas teaches that the blessed see God but never comprehend him, and he puts it with a saying he takes from Augustine:

\begin{quote}
It is impossible for any created intellect to comprehend God; yet ``for the mind to attain to God in some degree is great beatitude,'' as Augustine says.\endnote{Aquinas, \emph{Summa theologiae} I, q.~12, a.~7, trans.\ English Dominican Province. Aquinas gives the Augustinian locus as \emph{De Verb.\ Dom.}, Serm.\ xxxviii.}
\end{quote}

His reason has two parts. God is infinitely knowable, and the light of glory, being created, is finite, so that no created intellect can know God as far as he can be known: ``Hence it is impossible that it should comprehend God.''\endnote{Aquinas, \emph{Summa theologiae} I, q.~12, a.~7, trans.\ English Dominican Province.} The saints see the whole God, for God has no parts; but they do not see him wholly, as he sees himself.

The psalmist prays in words that fit this teaching: ``My heart hath said to thee: My face hath sought thee: thy face, O Lord, will I still seek.''\endnote{Ps 26:8, Douay--Rheims (Challoner).} The seeking that begins in faith ends in sight, and the sight never exhausts what it sees.
